% Cote 119: offprints, reports, courses, CV and letters (typed and printed part): archive cote 119, site page 72 (PDF page 73).
% offprint. Draft compared: none.
% AI-transcribed by Claude Opus 5.5 (claude-opus-5-5, adaptive thinking, effort high), Message Batches API, batch r119_full (msgbatch_01DwZNm5Y7YsKBWyAN4sFs8w), 2026-11-13; not yet checked by eye.
\scanpage{119/72}
\mspage{( 4 )}
(i) G est de type fini sur $k$ si et seulement si la $\otimes$-catégorie $\mathrm{Rep}_0(\mathrm{G})$
possède un $\otimes$-générateur.

(ii) L'action de G dans V est fidèle si et seulement si $\mathrm{V}\oplus\check{\mathrm{V}}$ est un $\otimes$-géné-%
rateur de $\mathrm{Rep}_0(\mathrm{G})$, ou encore si et seulement si tout objet de $\mathrm{Rep}_0(\mathrm{G})$ est
quotient d'un sous-objet d'un objet de la forme $\mathrm{P}(\mathrm{V},\check{\mathrm{V}})$, $\mathrm{P}\in\mathbf{N}[t_1, t_2]$.

(iii) $u$ est un épimorphisme de faisceaux fpqc si et seulement si $u$ est un
morphisme fidèlement plat, ou encore si et seulement si
\[ \mathrm{Rep}_0(u) : \mathrm{Rep}_0(\mathrm{G}') \to \mathrm{Rep}_0(\mathrm{G}) \]
est pleinement fidèle et identifie $\mathrm{Rep}_0(\mathrm{G}')$ à une sous-catégorie de $\mathrm{Rep}_0(\mathrm{G})$
stable par sous-objets.

(iv) $u$ est un monomorphisme si et seulement si $u$ est une immersion fermée,
ou encore si et seulement si tout objet de $\mathrm{Rep}_0(\mathrm{G})$ est quotient d'un sous-%
objet d'un objet de la forme $\mathrm{Rep}_0(u)(\mathrm{V}')$, où $\mathrm{V}'\in\mathrm{ob}[\mathrm{Rep}_0(\mathrm{G}')]$.

(*) Séance du 11 janvier 1971.
(${}^1$) M. Artin, A. Grothendieck et J.-L. Verdier, SGA 4, I. H. E. S.
(${}^2$) J. Benabou, Comptes rendus, 256, 1963, p. 1887. Le signe ACU est pour Associatif,
Commutatif, Unifère, et on assume que les contraintes ACU sont compatibles entre elles.
Enfin, pour nous, si F est un $\otimes$-foncteur, la flèche $\mathrm{F(X)}\otimes\mathrm{F(Y)}\to\mathrm{F(X}\otimes\mathrm{Y)}$ est un
isomorphisme.
(${}^3$) P. Cartier, Comptes rendus, 242, 1956, p. 322.
(${}^4$) M. Demazure et A. Grothendieck, SGA 3, Lecture Notes in Mathematics n${}^{\mathrm{os}}$ 151,
152 et 153, Springer-Verlag.
(${}^5$) J. Giraud, Cohomologie non abélienne, Notes miméographiées, Columbia University.
(${}^6$) A. Grothendieck, SGA 1, I. H. E. S.
(${}^7$) On dit qu'une $\otimes$-catégorie ACU $\mathcal{C}$ possède des objets Hom si, pour deux objets
quelconques X, Y de $\mathcal{C}$, le foncteur $\mathrm{Z}\mapsto\mathrm{Hom}(\mathrm{Z}\otimes\mathrm{X},\mathrm{Y})$ est représentable.

(Institut des Hautes Études
Scientifiques,
35, route de Chartres,
91-Bures-sur-Yvette,
Essonne.)

183186. --- Imp. Gauthier-Villars. --- 55, Quai des Grands-Augustins, Paris (6${}^{\mathrm{e}}$).
Imprimé en France.

\margin{72}
